% !TEX root = ../Main.tex
\chapter{计数原理}

概率的基础是\textbf{计数}——事件发生的"有利情况数"与"所有情况数"之比。两条基本原理串联了所有计数问题：\textbf{分类加法} 与 \textbf{分步乘法}。

\section{分类加法计数原理}

\state{分类加法计数原理}{
    完成一件事有 $n$ 类办法，第 $i$ 类办法中有 $m_i$ 种不同的方法，则完成这件事共有
    \[
    m_1 + m_2 + \cdots + m_n
    \]
    种不同的方法。

    \textbf{关键词：} "类"——不同类之间\textbf{互斥}，完成这件事只需选用\textbf{任一类}中的\textbf{任一种}方法。
}

\section{分步乘法计数原理}

\state{分步乘法计数原理}{
    完成一件事需要 $n$ 个步骤，第 $i$ 步有 $m_i$ 种不同的做法，则完成这件事共有
    \[
    m_1 \cdot m_2 \cdots m_n
    \]
    种不同的方法。

    \textbf{关键词：} "步"——每一步都\textbf{必须做}，每一步内部独立地有多种选择。
}

\ex[从"加"到"乘"的组合使用]{
    某同学从学校到家有 $3$ 条直达路线，也可选择先到图书馆再回家——学校到图书馆有 $2$ 条路，图书馆到家有 $4$ 条路。问从学校到家共有多少种走法？
}

\sol{
    \textbf{先分类：}
    \begin{itemize}
        \item 类一（直达）：$3$ 种走法；
        \item 类二（经过图书馆）：需\textbf{分两步}——先选学校→图书馆（$2$ 种），再选图书馆→家（$4$ 种），由乘法原理此类共 $2 \times 4 = 8$ 种。
    \end{itemize}
    \textbf{两类互斥，由加法}：$3 + 8 = 11$ 种。
}

\rmkb{
    \textbf{加与乘的辨别}：
    \begin{itemize}
        \item \textbf{加法（"分类"）}：完成任务有\textbf{不同方案}，每方案独立可达目标——"做其一即可"；
        \item \textbf{乘法（"分步"）}：完成任务需\textbf{多个步骤}，缺一不可——"都要做"。
    \end{itemize}
    实际问题中两者常交织——\textbf{先分类再分步}（如本例）或\textbf{先分步再分类}。识别清楚"类"与"步"的结构是计数的第一步。
}

\section{排列与组合}

\state{排列与组合}{
    从 $n$ 个不同元素中取出 $k$ 个：
    \begin{itemize}
        \item \textbf{排列}（关心顺序）：$\displaystyle A_n^k = \frac{n!}{(n-k)!} = n(n-1)\cdots(n-k+1)$；
        \item \textbf{组合}（不关心顺序）：$\displaystyle C_n^k = \frac{n!}{k!(n-k)!} = \frac{A_n^k}{k!}$。
    \end{itemize}
    关系：$A_n^k = C_n^k \cdot k!$——排列 = 先选后排。
}
